\documentclass[12pt,a4paper]{report}

% =========================================================================
%                         1. PACKAGES & SETUP
% =========================================================================
\usepackage[utf8]{inputenc}
\usepackage[T1]{fontenc}
\usepackage[english]{babel} % Change to 'french' if needed
\usepackage{geometry}       % For margins
\usepackage{graphicx}       % For images
\usepackage{xcolor}         % For colors
\usepackage{titlesec}       % For custom chapter titles
\usepackage{eso-pic}        % For background images
\usepackage{float}          % For [H] figure placement
\usepackage{subcaption}     % For subfigures
\usepackage{hyperref}       % For clickable links
\usepackage{booktabs}       % For nicer tables

% --- Code Highlighting ---
% NOTE: 'minted' requires Python and the -shell-escape flag. 
% If you don't have this set up, keep this commented and use 'listings' or 'verbatim'.
%\usepackage{minted} 
%\definecolor{bg}{rgb}{0.95,0.95,0.95} 

% =========================================================================
%                  2. CUSTOM COMMANDS (Replacing the .cls file)
% =========================================================================

% --- Variables for the Title Page ---
\newcommand{\titre}[1]{\def\thetitle{#1}}
\newcommand{\UE}[1]{\def\theue{#1}}
\newcommand{\sujet}[1]{\def\thesujet{#1}}
\newcommand{\enseignant}[1]{\def\theenseignant{#1}}
\newcommand{\eleves}[1]{\def\theeleves{#1}}
\newcommand{\jury}[1]{\def\thejury{#1}}

% --- Margin Setup Command ---
\newcommand{\fairemarges}{
    \geometry{top=2.5cm, bottom=2.5cm, left=2.5cm, right=2.5cm}
}

% --- Title Page Layout ---
\newcommand{\fairepagedegarde}{
    \begin{titlepage}
        \thispagestyle{empty}
        \centering
        % LOGO PLACEHOLDER
        % \includegraphics[width=0.4\textwidth]{images/logo.png} 
        \vspace*{1cm}
        
        {\scshape\LARGE \theue \par}
        \vspace{0.5cm}
        {\scshape\Large \thesujet \par}
        
        \vspace{2.5cm}
        
        \rule{\linewidth}{0.5mm} \\[0.4cm]
        {\huge\bfseries \thetitle \par}
        \rule{\linewidth}{0.5mm} \\[1.5cm]
        
        \vspace{1cm}
        
        \begin{minipage}{0.45\textwidth}
            \begin{flushleft} \large
                \emph{Author:}\\
                \theeleves
            \end{flushleft}
        \end{minipage}
        \hfill
        \begin{minipage}{0.45\textwidth}
            \begin{flushright} \large
                \emph{Supervisors:} \\
                \theenseignant
            \end{flushright}
        \end{minipage}

        \vspace{2cm}
        
        \ifx\thejury\empty \else
            \vfill
            {\large \textbf{Jury:} \\ \thejury}
        \fi

        \vfill
        {\large \today\par} 
    \end{titlepage}
}

% =========================================================================
%                  3. CHAPTER STYLING & BACKGROUND
% =========================================================================

% --- Background Image Command ---
\newcommand{\chapterbackground}{%
  \AddToShipoutPictureBG*{%
    \AtPageLowerLeft{%
      % UNCOMMENT THE LINE BELOW ONLY IF YOU HAVE THE IMAGE
      % \includegraphics[width=\paperwidth,height=\paperheight,opacity=0.1]{images/backgroundcap.png}%
    }%
  }%
}

% --- Chapter Title Formatting ---
\titleformat{\chapter}[display]
  {\normalfont\bfseries\LARGE}
  {\mdseries\fontsize{60pt}{72pt}\selectfont\chaptertitlename\ \thechapter}
  {25pt}
  {\chapterbackground}

\titleformat{name=\chapter,numberless}
  {\normalfont\bfseries\LARGE} 
  {}{0pt}
  {\chapterbackground}


% =========================================================================
%                             4. DOCUMENT START
% =========================================================================
\begin{document}

% --- Define Report Info Here ---
\titre{Project Title Goes Here}
\UE{Industrial Engineering \& Data Science}
\sujet{End of Year Project}
\enseignant{Prof. Supervisor 1 \\ Prof. Supervisor 2}
\eleves{Student Name}
\jury{Prof. Jury Member}

% --- Initialization ---
\fairepagedegarde
\pagenumbering{roman} 
\setcounter{page}{1}
\fairemarges 

% =========================================================================
%                          PRELIMINARY PAGES
% =========================================================================

% --- Dedication ---
\cleardoublepage
\thispagestyle{empty}
\vspace*{4cm}
\begin{center} \textbf{\Huge Dedication} \end{center}
\vspace{2cm}
\begin{center} \Large
To my family, for their support...
\end{center}

% --- Acknowledgements ---
\newpage
\thispagestyle{empty}
\vspace*{4cm}
\begin{center} \textbf{\Huge Acknowledgements} \end{center}
\vspace{2cm}
\begin{center} \Large
I would like to thank my supervisors...
\end{center}

% --- Abstract ---
\chapter*{Abstract}
\addcontentsline{toc}{chapter}{Abstract}
This report details the development of... [Your abstract here]

\newline \newline \hline \vspace{0.5cm}
{\large \textbf{Keywords}: } Python, Automation, Data Science, Quality Control.
\vspace{0.5cm} \hline
\newpage

% --- List of Abbreviations ---
\chapter*{List of Abbreviations}
\addcontentsline{toc}{chapter}{List of Abbreviations}
\begin{tabular}{ll}
    \textbf{API} & Application Programming Interface \\
    \textbf{GUI} & Graphical User Interface \\
    \textbf{PCP} & Part Control Plan \\
\end{tabular}
\newpage

% --- Glossary ---
\chapter*{Glossary}
\addcontentsline{toc}{chapter}{Glossary}
\begin{description}
    \item[Term 1] Definition of term 1.
    \item[Term 2] Definition of term 2.
\end{description}
\newpage

% --- TOC ---
\tabledematieres % French command mapping isn't standard, using standard below
\tableofcontents
\listoffigures
\listoftables
\newpage

% =========================================================================
%                              MAIN CONTENT
% =========================================================================
\pagenumbering{arabic} 
\setcounter{page}{1}

% --- Introduction ---
\chapter*{General Introduction}
\addcontentsline{toc}{chapter}{General Introduction}
Introduction text goes here...

\chapter{Project Context}
\section{Host Organization}
Description of the company...
\section{Problem Statement}
Description of the problem...

\chapter{Technical Implementation}
\section{Architecture}
The system follows a 3-tier architecture...

\section{Code Implementation}
% Example of code. 
% NOTE: If you uncommented 'minted' in the preamble, you can use \begin{minted}{python}
\begin{verbatim}
def analyze_data(data):
    # This is a code block placeholder
    return data.process()
\end{verbatim}

\chapter{Results and Discussion}
\section{Performance}
The tool reduced time by 99\%...

\begin{table}[H]
    \centering
    \caption{Time Comparison}
    \begin{tabular}{lcc}
        \toprule
        \textbf{Task} & \textbf{Manual} & \textbf{Automated} \\
        \midrule
        Analysis & 2 hours & 5 mins \\
        Reporting & 1 hour & Instant \\
        \bottomrule
    \end{tabular}
\end{table}

\chapter*{General Conclusion}
\addcontentsline{toc}{chapter}{General Conclusion}
We successfully achieved the objectives...

% =========================================================================
%                             BIBLIOGRAPHY
% =========================================================================
% You can use a .bib file, or write references manually here for a single file solution
\begin{thebibliography}{9}
    \bibitem{zheng2020smart}
    Zheng, P. et al. (2020). \textit{Smart manufacturing systems for Industry 4.0}.
    
    \bibitem{pcpcontext}
    Automotive Industry Action Group. (2008). \textit{Advanced Product Quality Planning}.
\end{thebibliography}
\addcontentsline{toc}{chapter}{Bibliography}

% =========================================================================
%                               ANNEXES
% =========================================================================
\chapter*{Annexes}
\addcontentsline{toc}{chapter}{Annexes}
\markboth{Annexes}{Annexes}

\section*{Annex A: Screenshots}
\begin{figure}[H]
    \centering
    % \includegraphics[width=0.8\textwidth]{images/screenshot.png}
    \caption{Screenshot placeholder}
\end{figure}

\end{document}